% !TEX root = ../main.tex
\section{Identity Escrow Protocol}\label{sec:escrow}

Zero-knowledge credential systems face a fundamental tension:
privacy requires that presentations reveal nothing beyond the truth
value of the requested predicates, yet the rule of law demands that,
under judicial authorisation, the identity behind a fraudulent or
harmful transaction can be recovered.
This section presents the \textsc{zk-eidas} identity escrow protocol,
which resolves this tension by encrypting identifying fields
\emph{per presentation} under a designated escrow authority's
post-quantum public key, with an in-circuit binding that
prevents substitution of the encrypted payload.

%% ───────────────────────────────────────────────────────────────────
\subsection{Motivation and Regulatory Requirement}\label{sec:escrow-motivation}

Article~5a of the eIDAS~2.0 Regulation requires that member states
ensure the possibility of identifying natural persons behind
electronic transactions when mandated by a court order or other
legitimate judicial authorisation.  An anonymous credential system
that provides \emph{no} mechanism for accountability creates a legal
vacuum: it cannot be deployed in any eIDAS-regulated context, because
no lawful process exists to trace a fraudulent presentation back to
its holder.

\paragraph{Existing approaches and their limitations.}
The literature offers several mechanisms for revocable anonymity, each
with significant drawbacks in the eIDAS setting:

\begin{itemize}
  \item \textbf{Group signatures}~\cite{camenisch2001}.
    A group manager can open \emph{any} signature produced by any group
    member, creating a single point of failure.  If the group manager is
    compromised, every past and future presentation is deanonymisable.
    Moreover, the opening power is unconditional --- it is not
    scoped to a court order or a specific presentation.

  \item \textbf{Trusted third parties (online escrow).}
    Some designs require the holder to contact an online escrow agent
    at each presentation to obtain a blinded token.  This introduces an
    availability risk (the escrow agent becomes a service bottleneck)
    and a correlation risk (the agent learns the timing and frequency
    of presentations, even if it does not learn the content).

  \item \textbf{Revocable anonymity via accumulators.}
    Accumulator-based revocation schemes~\cite{camenisch2001} are
    designed for credential \emph{revocation} (removing a holder from
    the valid set), not for per-presentation accountability.  They do
    not provide a mechanism to recover the identity behind a specific
    transaction; they can only exclude a known holder from future
    transactions.
\end{itemize}

\paragraph{Our approach.}
We adopt a \emph{per-presentation encryption} strategy.  At each
proof generation, the holder encrypts a configurable set of
identifying fields (e.g., name and date of birth) under the escrow
authority's ML-KEM-768 encapsulation key~\cite{fips203}.  The
resulting ciphertext is bound to the zero-knowledge proof via an
in-circuit SHA-256 digest, so that the verifier can later confirm ---
after a court-ordered decryption --- that the recovered identity
matches the proof.  Crucially:
\begin{enumerate}
  \item Each presentation uses a \emph{fresh} ML-KEM encapsulation,
    so ciphertexts from different presentations are unlinkable
    (Theorem~\ref{thm:unlinkability}).
  \item The escrow authority holds only a decapsulation key; it cannot
    decrypt without possessing the per-presentation ciphertext, and it
    learns nothing from observing the proof itself.
  \item No online interaction with the authority is required at
    presentation time.
\end{enumerate}

%% ───────────────────────────────────────────────────────────────────
\subsection{Construction}\label{sec:escrow-construction}

We describe the escrow protocol as four procedures:
$\mathsf{KeyGen}$ (authority setup),
$\mathsf{Escrow}$ (holder encryption during proof generation),
$\mathsf{Bind}$ (in-circuit digest computation), and
$\mathsf{Recover}$ (authority decryption upon court order).

\medskip

\begin{algorithm}[t]
  \caption{$\mathcal{E}.\mathsf{KeyGen}$: Escrow authority key generation}\label{alg:keygen}
  \begin{algorithmic}[1]
    \Ensure Encapsulation key $\mathit{ek}$; seed $s$
    \Statex
    \State $s \gets\!\!\!\!\$\ \{0,1\}^{512}$
      \Comment{64-byte random seed}
    \State $\mathit{dk} \gets \textsf{ML-KEM-768}.\mathsf{DeriveKey}(s)$
      \Comment{FIPS 203~\cite{fips203}}
    \State $\mathit{ek} \gets \mathit{dk}.\mathsf{EncapsulationKey}()$
      \Comment{1184 bytes}
    \State Publish $\mathit{ek}$
    \State Store $s$ securely
      \Comment{offline, hardware-protected}
    \State \Return $(\mathit{ek}, s)$
  \end{algorithmic}
\end{algorithm}

\begin{algorithm}[t]
  \caption{$\mathcal{H}.\mathsf{Escrow}$: Per-presentation field encryption (holder side)}\label{alg:escrow}
  \begin{algorithmic}[1]
    \Require Field values $\{f_i\}_{i=0}^{n-1}$ ($n \leq 8$, each a byte string); encapsulation key $\mathit{ek}$
    \Ensure Envelope $\mathcal{E}\mathit{nv} = (\mathit{cts}[], \mathit{tags}[], \mathit{ek\_blob}, \mathit{names}[])$
    \Statex
    \State $K \gets\!\!\!\!\$\ \{0,1\}^{248}$
      \Comment{31-byte symmetric key, zero-padded to 32 bytes}
    \For{$i = 0, \ldots, n-1$}
      \State $\mathit{nonce}_i \gets 0^{64} \,\|\, i$
        \Comment{12-byte nonce, counter mode}
      \State $(\mathit{ct}_i, \mathit{tag}_i) \gets \textsf{AES-256-GCM}.\mathsf{Enc}(K, \mathit{nonce}_i, f_i)$
    \EndFor
    \State $(c, \mathit{ss}) \gets \textsf{ML-KEM-768}.\mathsf{Encaps}(\mathit{ek})$
      \Comment{$c$: 1088 bytes}
    \State $\mathit{mask} \gets \mathsf{SHA\textnormal{-}256}(\mathit{ss})$
      \Comment{32-byte mask}
    \State $\mathit{enc}_K \gets K \oplus \mathit{mask}$
      \Comment{32 bytes}
    \State \Return $\bigl(\{\mathit{ct}_i\}, \{\mathit{tag}_i\}, c \,\|\, \mathit{enc}_K, \{\mathit{name}_i\}\bigr)$
  \end{algorithmic}
\end{algorithm}

\begin{algorithm}[t]
  \caption{$\mathsf{Bind}$: In-circuit escrow digest (inside Longfellow hash sub-circuit)}\label{alg:bind}
  \begin{algorithmic}[1]
    \Require Escrow field values $\{e_j\}_{j=0}^{7}$, each $\in \{0,1\}^{256}$ (32-byte zero-padded claim values)
    \Ensure Public output $\delta$ (escrow digest)
    \Statex
    \State $\delta \gets \mathsf{SHA\textnormal{-}256}(e_0 \,\|\, e_1 \,\|\, \cdots \,\|\, e_7)$
      \Comment{256-byte input, single hash}
    \State \textbf{output} $\delta$ as public
  \end{algorithmic}
\end{algorithm}

Algorithm~\ref{alg:bind} is executed within the same Longfellow circuit~\cite{longfellow} that verifies the credential (Algorithm~\ref{alg:prove}).  The escrow digest~$\delta$ appears as a public output alongside the nullifier~$\eta$ and binding hash~$\beta$.  Because the circuit takes the raw claim values as private input and hashes them deterministically, the digest commits to the \emph{actual} credential fields --- not to values chosen by the holder after the fact.

The verifier receives the proof (containing~$\delta$) and the envelope produced by Algorithm~\ref{alg:escrow}.  It stores both.  The binding is verified \emph{only after court-ordered decryption}: the authority decrypts the envelope, pads each recovered field to 32~bytes, computes $\mathsf{SHA\textnormal{-}256}(e_0 \| \cdots \| e_7)$, and checks equality with~$\delta$.  A mismatch indicates tampering.

\begin{algorithm}[t]
  \caption{$\mathcal{E}.\mathsf{Recover}$: Identity recovery upon court order}\label{alg:recover}
  \begin{algorithmic}[1]
    \Require Envelope $\mathcal{E}\mathit{nv}$; authority seed $s$
    \Ensure Plaintext fields $\{f_i\}_{i=0}^{n-1}$
    \Statex
    \State $\mathit{dk} \gets \textsf{ML-KEM-768}.\mathsf{DeriveKey}(s)$
    \State Parse $c \,\|\, \mathit{enc}_K$ from $\mathcal{E}\mathit{nv}.\mathit{ek\_blob}$
    \State $\mathit{ss} \gets \textsf{ML-KEM-768}.\mathsf{Decaps}(\mathit{dk}, c)$
    \State $\mathit{mask} \gets \mathsf{SHA\textnormal{-}256}(\mathit{ss})$
    \State $K \gets \mathit{enc}_K \oplus \mathit{mask}$
    \For{$i = 0, \ldots, n-1$}
      \State $\mathit{nonce}_i \gets 0^{64} \,\|\, i$
      \State $f_i \gets \textsf{AES-256-GCM}.\mathsf{Dec}(K, \mathit{nonce}_i, \mathit{ct}_i \,\|\, \mathit{tag}_i)$
    \EndFor
    \State \Return $\{f_i\}_{i=0}^{n-1}$
  \end{algorithmic}
\end{algorithm}

\paragraph{Concrete sizes.}
For a typical escrow envelope containing two fields (family name and date of birth), the overhead is dominated by the ML-KEM-768 ciphertext: $1088 + 32 = 1120$~bytes for the key blob, plus approximately $2 \times (32 + 16) = 96$~bytes for the AES-256-GCM ciphertexts and tags, totalling roughly 1.2~KB.  This is appended to the proof, which itself is approximately 365~KB (see Section~\ref{sec:benchmarks}).

\begin{figure}[t]
\centering
% !TEX root = ../main.tex
% Figure: Identity escrow protocol sequence diagram
\begin{tikzpicture}[
  node distance=0 and 0,
  arrow/.style={-{Stealth[length=5pt]}, thick},
  dasharrow/.style={-{Stealth[length=5pt]}, thick, dashed},
  every node/.style={font=\small}
]

%% ── Column x-positions ───────────────────────────────────────────────
\def\xH{0}      % Holder
\def\xV{5.5}    % Verifier
\def\xE{11}     % Escrow Authority

%% ── Party headers ────────────────────────────────────────────────────
\node[align=center] (LH) at (\xH, 0) {\textbf{Holder} ($\mathcal{H}$)};
\node[align=center] (LV) at (\xV, 0) {\textbf{Verifier} ($\mathcal{V}$)};
\node[align=center] (LE) at (\xE, 0) {\textbf{Escrow Authority} ($\mathcal{E}$)};

%% ── Vertical timelines ───────────────────────────────────────────────
\draw[thick] (\xH,  -0.3) -- (\xH, -7.8);
\draw[thick] (\xV,  -0.3) -- (\xV, -7.8);
\draw[thick] (\xE,  -0.3) -- (\xE, -7.8);

%% ── Step 1: Holder → Verifier: ZK proof + escrow envelope ────────────
\draw[arrow] (\xH, -1.0) -- (\xV, -1.0)
    node[midway, above, font=\small] {ZK proof $+$ escrow envelope};

%% ── Step 2: Verifier internal: verify proof ──────────────────────────
\draw[thick] (\xV-0.1, -1.7) rectangle (\xV+0.1, -2.5);
\node[right=0.15cm, align=left, font=\small] at (\xV+0.1, -2.1)
    {verify proof\\accept / reject};

%% ── Separator: normal flow ───────────────────────────────────────────
\draw[dotted, thick] (-0.7, -3.2) -- (\xE+0.8, -3.2);
\node[font=\scriptsize\itshape, right] at (-0.7, -3.15) {};

%% ── Court-order dashed box ───────────────────────────────────────────
\draw[dashed, rounded corners=4pt, thick]
    (-0.5, -3.5) rectangle (\xE+0.6, -6.2);
\node[font=\small\itshape, anchor=north west] at (-0.5, -3.5)
    {[court order]};

%% ── Step 3: Verifier → Escrow Authority: envelope ────────────────────
\draw[arrow] (\xV, -4.2) -- (\xE, -4.2)
    node[midway, above, font=\small] {envelope};

%% ── Step 4: Escrow Authority internal: ML-KEM + AES-GCM decrypt ──────
\draw[thick] (\xE-0.1, -4.6) rectangle (\xE+0.1, -5.8);
\node[right=0.15cm, align=left, font=\small] at (\xE+0.1, -5.2)
    {ML-KEM-768 decaps\\AES-256-GCM decrypt\\bind check ($\delta$)};

%% ── Step 5: Escrow Authority → Verifier: recovered identity ──────────
\draw[arrow] (\xE, -6.5) -- (\xV, -6.5)
    node[midway, above, font=\small] {identity fields $\{f_i\}$};

%% ── Step 6: Verifier → Holder (optional notification, dashed) ────────
\draw[dasharrow] (\xV, -7.3) -- (\xH, -7.3)
    node[midway, above, font=\small] {\textit{(optional notification)}};

\end{tikzpicture}

\caption{Identity escrow protocol flow. The escrow envelope is transmitted alongside the proof. Decryption occurs only upon court order.}
\label{fig:escrow-protocol}
\end{figure}

%% ───────────────────────────────────────────────────────────────────
\subsection{Security Analysis}\label{sec:escrow-security}

We state four security theorems for the escrow protocol.  We provide proof sketches here; full game-based reductions are deferred to the extended version of this paper.

\begin{theorem}[Escrow Correctness]\label{thm:correctness}
  If the holder $\mathcal{H}$ follows the protocol honestly --- executing
  $\mathsf{Escrow}$ (Algorithm~\ref{alg:escrow}) with a valid
  encapsulation key $\mathit{ek}$ and the same field values used as
  private inputs to the Longfellow circuit --- then the escrow authority
  $\mathcal{E}$ recovers the correct identity fields via
  $\mathsf{Recover}$ (Algorithm~\ref{alg:recover}), and the recovered
  fields satisfy $\mathsf{SHA\textnormal{-}256}(e_0 \| \cdots \| e_7) = \delta$.
\end{theorem}

\begin{proof}[Proof sketch]
  Correctness follows from two standard properties.
  First, ML-KEM-768 \emph{decapsulation correctness} (FIPS~203, Theorem~1~\cite{fips203}):
  for a correctly generated key pair $(\mathit{dk}, \mathit{ek})$ and
  ciphertext $(c, \mathit{ss}) \gets \mathsf{Encaps}(\mathit{ek})$,
  the decapsulation $\mathsf{Decaps}(\mathit{dk}, c)$ returns
  the same shared secret~$\mathit{ss}$ with probability~$1 - 2^{-138}$.
  Since $\mathit{mask} = \mathsf{SHA\textnormal{-}256}(\mathit{ss})$ is deterministic,
  the authority recovers the same mask and hence the same
  key~$K = \mathit{enc}_K \oplus \mathit{mask}$.
  Second, AES-256-GCM \emph{decryption correctness}: given the correct
  key and nonce, decryption inverts encryption exactly.
  The deterministic nonce derivation ($0^{64} \| i$) ensures each
  field is decrypted with the correct nonce.
  The escrow digest check succeeds because the circuit and the
  encryption operate on identical field values.
\end{proof}

\begin{theorem}[Non-Frameability]\label{thm:non-frameability}
  No probabilistic polynomial-time adversary $\mathcal{A}$ (even one
  controlling the escrow authority $\mathcal{E}^*$ and the verifier
  $\mathcal{V}^*$) can produce a pair $(\pi, \mathcal{E}\mathit{nv})$
  such that (a)~the proof $\pi$ is accepted by an honest verifier with
  escrow digest~$\delta$, and (b)~$\mathcal{E}\mathit{nv}$ decrypts
  to fields belonging to an honest holder $\mathcal{H}'$ who did not
  participate, except with negligible probability.  Formally: if
  $\mathcal{A}$ frames with probability~$\varepsilon$, there exists a
  reduction to a SHA-256 second-preimage finder with advantage
  $\geq \varepsilon$.
\end{theorem}

\begin{proof}[Proof sketch]
  The escrow digest $\delta$ is computed \emph{inside} the Longfellow
  circuit over the actual claim values that satisfy the credential
  verification constraints.  To frame a holder~$\mathcal{H}'$, the
  adversary must produce an envelope whose decrypted fields
  $\{f'_i\}$ hash to the same~$\delta$ as the genuine fields
  $\{e_j\}$, i.e., must find a second preimage for SHA-256 on a
  256-byte input.  Given such an adversary, we construct a
  second-preimage finder: run $\mathcal{A}$, extract the framed
  fields, and output them as the second preimage.  By the
  second-preimage resistance of SHA-256, this succeeds with
  negligible probability.
\end{proof}

\begin{theorem}[Unlinkability Preservation]\label{thm:unlinkability}
  The escrow ciphertexts do not help a probabilistic polynomial-time
  adversary link two presentations from the same holder.  That is,
  the escrow mechanism does not degrade the unlinkability property
  of Definition~\ref{def:unlinkability}.
\end{theorem}

\begin{proof}[Proof sketch]
  Each invocation of $\mathsf{Escrow}$ (Algorithm~\ref{alg:escrow})
  draws a fresh random coin for $\textsf{ML-KEM-768}.\mathsf{Encaps}$,
  producing a fresh ciphertext~$c$ and a fresh shared secret~$\mathit{ss}$.
  From the IND-CCA2 security of ML-KEM-768 under the Module-LWE
  assumption~\cite{fips203}, the shared secret is computationally
  indistinguishable from a uniform random value.  Consequently,
  $\mathit{mask} = \mathsf{SHA\textnormal{-}256}(\mathit{ss})$ is
  pseudorandom, $\mathit{enc}_K = K \oplus \mathit{mask}$ is
  one-time-pad--encrypted, and the AES-256-GCM ciphertexts under
  distinct keys are IND-CCA2 secure.  Even if two presentations
  encrypt the \emph{same} identity fields, the resulting envelopes
  are computationally indistinguishable from encryptions of unrelated
  plaintexts, so no distinguisher can link them.
\end{proof}

\begin{theorem}[Post-Quantum Escrow Security]\label{thm:pq-security}
  The escrow envelope is IND-CCA2 secure against quantum adversaries
  at NIST security Level~3.
\end{theorem}

\begin{proof}[Proof sketch]
  The construction follows the KEM/DEM paradigm.
  The KEM component is ML-KEM-768, which achieves NIST security
  Level~3 ($\geq$ AES-192 equivalent) under the Module-LWE
  hardness assumption~\cite{fips203}.
  The DEM component is AES-256-GCM, which retains at least 128-bit
  security against quantum adversaries under Grover's bound
  (i.e., the effective key strength halves from 256 to 128~bits,
  still exceeding the 128-bit target of NIST Level~3).
  By the KEM/DEM composition theorem~\cite{cramer2003design}, if
  the KEM is IND-CCA2 and the DEM is one-time IND-CCA2, the
  composed scheme is IND-CCA2.  Here, each symmetric key~$K$ is
  used exactly once (fresh per presentation), so the one-time
  requirement is satisfied.  The overall scheme therefore inherits
  the post-quantum security of the Module-LWE assumption.
\end{proof}
